\chapter{Validatie tijdens de test in Spa}
\label{chap:spa}

\pending{De praktijktest in Spa heeft nog niet plaatsgevonden. Dit hoofdstuk beschrijft het validatieplan, niet de resultaten. Alle meetwaarden, observaties en conclusies moeten na de test worden ingevuld.}

\section{Doel van de test}
De test in Spa moet aantonen of de applicatie buiten demo- en historische data bruikbaar blijft gedurende een echte operationele sessie. De nadruk ligt niet alleen op correcte cijfers, maar ook op installatie, stabiliteit, begrijpelijkheid en de snelheid waarmee de race-engineer informatie kan interpreteren.

\section{Voorbereiding}
Voor de sessie moeten timingprovider, URL, gevolgde wagens, sessiemodus en een lege opslagmap worden ingesteld. Voor racegebruik moeten raceduur, verplichte stops, gekozen Spa-pitconfiguratie, reguliere trajectafstand en FCY-snelheid worden gecontroleerd met het actuele reglement. Referentietijden worden pas ingevoerd wanneer de opdrachtgever de juiste waarden bevestigt.

\pending{Noteer hier de gebruikte appversie, laptop en besturingssysteem, timing-URL, Spa-layout, betrokken autonummers, sessietype, reglementaire pitstopregels en referentietijden.}

\section{Te valideren scenario's}
Minstens de volgende punten moeten tijdens of onmiddellijk na de test worden gecontroleerd:

\begin{enumerate}
  \item Komt iedere voltooide ronde precies eenmaal in de historiek terecht?
  \item Komen pilootwissels en actuele sectoren tijdig op het juiste dashboard?
  \item Worden inlaps, outlaps en neutralisatieronden correct opgeslagen en uitgesloten?
  \item Sluit de predicted lap time na sector 1 en 2 voldoende aan bij de uiteindelijke ronde?
  \item Reageren referentiekleuren en delta's op de geconfigureerde regels?
  \item Komen verschillen met voorliggende en achterliggende klassewagens overeen met de officiële timing?
  \item Is de pitstopprojectie voor green flag en FCY plausibel vóór en na een echte stop?
  \item Blijven cooldown, verplichte-stopteller en herstartstatus correct?
  \item Blijven maximaal drie dashboards en het grafiekvenster synchroon?
  \item Kan de app na sluiten of een gesimuleerde crash met dezelfde map hervatten?
\end{enumerate}

\section{Meetplan}
Voor elke voorspelde ronde worden voorspelling na sector 1, voorspelling na sector 2 en werkelijke rondetijd geregistreerd. Daaruit kunnen absolute fouten worden berekend. Voor pitstops worden voorspelde klassepositie, voorspelde gap en werkelijke toestand na stabilisatie genoteerd. Bij afwijkingen moet ook de vlagstatus en kwaliteit van de timingfeed worden bewaard.

\begin{table}[ht]
\centering
\caption{In te vullen meettabel voor rondevoorspellingen in Spa.}
\begin{tabular}{rrrrr}
\toprule
Ronde & Piloot & Voorspelling S1 & Voorspelling S2 & Werkelijk \\
\midrule
-- & -- & -- & -- & -- \\
-- & -- & -- & -- & -- \\
\bottomrule
\end{tabular}
\end{table}

\pending{Vul na de test gemiddelde absolute fout, grootste afwijking en observaties per weers- of vlagconditie in. Voeg geen geschatte waarden toe wanneer een meting ontbreekt.}

\section{Gebruikersfeedback}
Naast technische metingen moet de opdrachtgever beoordelen welke velden daadwerkelijk tijdens de sessie werden bekeken, welke waarschuwingen duidelijk waren en welke informatie te veel plaats innam. Ook moet worden gevraagd of de pitstopprojectie voldoende uitlegbaar is om er een beslissing op te baseren.

\pending{Voeg hier concrete feedback van de opdrachtgever toe, inclusief voorgestelde wijzigingen en hun prioriteit.}

\section{Resultaat en besluit}
\pending{Schrijf na de Spa-test een feitelijk besluit: welke onderdelen zijn gevalideerd, welke fouten zijn gevonden, welke fixes zijn uitgevoerd en welke risico's blijven open. Tot dan mag dit hoofdstuk niet als bewijs van praktijkvalidatie worden gebruikt.}

